\section{A stable class with four-particle clusters}
\label{sec:model}

We first construct a class of potentials for which thermodynamic limits
exist and the required density bounds are uniform in the ranges of the
attractive interactions.
The positive core permits small clusters of at most four particles without
an energy cost. Its quadratic penalty for larger local populations will
control every attractive term added later.

\subsection{The packing density and the core}

Let $M_L$ be the largest cardinality of a subset of
$\Lambda_L=[-L/2,L/2]^3$ whose distinct points have distance at least one.
The number is finite, and compactness gives a maximizing configuration.
There is a constant
\begin{equation}
 p=\lim_{L\to\infty}\frac{M_L}{L^3}\in(0,\infty),
 \qquad A=\frac{3}{2p}.
 \label{eq:packing-density}
\end{equation}
Here we need only the existence of $p$, not its numerical value. Disjoint
balls of radius $1/2$ around the packing points lie in the cube of side
$L+1$, giving an upper bound on $M_L/L^3$ for $L\ge1$. A unit lattice
gives a positive lower bound. For convergence, repeat a maximizing packing
in a cube of side $s$ with period $s+1$ in each coordinate. Points in
different copies are still separated by at least one, so
\[
 \liminf_{L\to\infty}\frac{M_L}{L^3}
 \ge \frac{M_s}{(s+1)^3}.
\]
Letting $s\to\infty$ along a sequence realizing the limsup proves
\eqref{eq:packing-density}.

Fix
\[
 \delta=10^{-3},\qquad b=\frac14.
\]
For a parameter $t>1$, put $E=t^2$ and $e_t=e^{-t}$. Define the radial
nonnegative core, with $r=|x|$, by
\begin{equation}
 \psi_t(x)=
 \begin{cases}
 +\infty,&r=0,\\
 E+(e_t/r)^{12},&0<r<e_t,\\
 0,&e_t\le r\le(1+\delta)e_t,\\
 E,&(1+\delta)e_t<r<1,\\
 0,&r\ge1.
 \end{cases}
 \label{eq:core}
\end{equation}
Thus particles can approach one another without paying the energy $E$
only when their separation lies in a narrow shell of radius $e_t$.
We allow $t$ to vary until the initial pressure configuration has been
constructed; thereafter one sufficiently large finite value will be fixed.

Choose a nonnegative radial function $w\in C^1(\R^3)$, supported in the
unit ball, strictly positive on the closed ball of radius $1/2$, and
normalized by $\int_{\R^3}w(x)\,\dd x=1$. For $R\ge1$ set
$w_R(x)=R^{-3}w(x/R)$. Our class consists of
\begin{equation}
 \varphi=\psi_t-S,
 \qquad S=t\sum_j\sigma_j w_{R_j},
 \qquad R_j\ge1,\quad \sigma_j\ge0,
 \quad \sum_j\sigma_j\le A+1.
 \label{eq:physical-class}
\end{equation}
The sum may be finite, empty, or countably infinite. Since
$\|w_R\|_\infty\le\|w\|_\infty$ and $\|w_R\|_1=1$, it converges
uniformly and in $L^1$ in the infinite case. In particular, $S$ is a
bounded continuous nonnegative function. Every potential in this class is
nonpositive at distances at least one. This sign will give a direct proof
of the canonical thermodynamic limit.

Partition $\R^3$ into half-open cubes $C_i$ of side $b$, indexed by
$i\in\mathbb Z^3$. For a finite configuration, let $n_i$ be the number
of its particles in $C_i$. In a finite cube we intersect these cells with
the cube; every resulting cell has volume at most $b^3$. Set
\begin{equation}
 J_{ij}=\sup_{x\in C_i,\,y\in C_j}S(x-y),
 \qquad \sum_jJ_{ij}\le M,
 \qquad M=c_b t(A+1),
 \label{eq:cell-row}
\end{equation}
where $c_b$ depends only on $b$ and $w$. To see the row bound, first use
one kernel $w_R$. Each nonzero cell supremum is at most
$\|w\|_\infty R^{-3}$, and there are at most $C_bR^3$ cells within
distance $R$ of a given cell because $R\ge1$. Sum this bound over the
coefficients $t\sigma_j$. Taking a supremum of a sum only decreases the
result relative to the sum of suprema. The same proof therefore applies
to infinitely many ranges and to cropped cells.

\begin{lemma}[The shell graph]
\label{lem:shell-graph}
No five points in $\R^3$ have all their mutual distances in
$[e_t,(1+\delta)e_t]$. Consequently, a graph whose vertices are points
of $\R^3$ and whose edges join pairs at these distances has at most
$3n^2/8$ edges on any set of $n$ vertices.
\end{lemma}

\begin{proof}
Suppose that five such points existed. Scale by $e_t^{-1}$ and form the
Gram matrix $G$ of the four vectors from one point to the other four.
Write $\varepsilon=2\delta+\delta^2$. The diagonal entries differ
from 1 by at most $\varepsilon$. By the polarization identity, each
off-diagonal entry differs from $1/2$ by at most $\varepsilon$.
The matrix $G_0$ with diagonal 1 and off-diagonal $1/2$ has smallest
eigenvalue $1/2$, whereas
\[
 \|G-G_0\|_{\mathrm{op}}\le4\varepsilon
 =0.008004<\frac12.
\]
Thus $G$ would have rank four, which is impossible for four vectors
in $\R^3$.

For completeness, we give the weighted proof of the required graph bound,
an instance of the Motzkin--Straus identity
\cite[Theorem~1]{MotzkinStraus1965}.
For a graph with no complete subgraph on five vertices, maximize
\[
 \sum_{\{i,j\}\text{ an edge}}a_i a_j
 \quad\text{over}\quad a_i\ge0,\qquad\sum_i a_i=1.
\]
Choose a maximizer with the smallest possible support. If two vertices
in its support are nonadjacent, transferring all their combined weight
to the one with the larger weighted neighbor sum cannot decrease the
objective. It decreases the support, a contradiction. The support is
therefore a clique with at most four vertices. At this maximizer the
objective is
\[
 \frac12\left(1-\sum_i a_i^2\right)\le\frac38.
\]
Applying this inequality to the weights $a_i=1/n$ gives the asserted
edge count. The case $n=0$ is immediate.
\end{proof}

\begin{proposition}[Uniform stability]
\label{prop:stability}
Put $A_0=E/8$ and $B_0=E/2$. Every potential in
\eqref{eq:physical-class} satisfies, for every finite configuration,
\begin{equation}
 \sum_{\ell<k}\psi_t(x_\ell-x_k)
 \ge A_0\sum_i n_i^2-B_0N,
 \qquad
 U_N^\varphi\ge (A_0-M/2)\sum_i n_i^2-B_0N.
 \label{eq:superstable}
\end{equation}
These inequalities remain valid if positive core interactions between
members of a partition into unions of whole cells are deleted, or if
any negative interactions are multiplied by factors in $[0,1]$.
\end{proposition}

\begin{proof}
The diameter of a cell is $\sqrt3/4<1$. Among the $n_i(n_i-1)/2$
pairs inside cell $i$, at most $3n_i^2/8$ lie in the shell by
Lemma~\ref{lem:shell-graph}. Every other pair has core energy at
least $E$. Thus the core energy inside this cell is at least
\[
 E\left(\frac{n_i(n_i-1)}2-\frac{3n_i^2}8\right)
 =A_0 n_i^2-B_0 n_i.
\]
All remaining core interactions are nonnegative. Also,
\[
 \sum_{\ell<k}S(x_\ell-x_k)
 \le\frac12\sum_{i,j}J_{ij}n_i n_j
 \le\frac M2\sum_i n_i^2,
\]
where the first bound includes harmless self-particle terms and the
second follows from symmetry, \eqref{eq:cell-row}, and
$2n_i n_j\le n_i^2+n_j^2$. This proves
\eqref{eq:superstable}. The permitted deletions preserve every core
interaction within a small cell. Reducing a negative term only
increases the energy, so the same argument proves the last assertion.
\end{proof}

From now on we require $t$ to be sufficiently large that
\begin{equation}
 2(1+\delta)e_t<1,
 \qquad A_0\ge8M,
 \qquad \kappa:=(A_0-M/2)b^3>10t(A+1).
 \label{eq:t-choice}
\end{equation}
These conditions hold for all sufficiently large $t$, since $A_0$ is
quadratic in $t$ and $M$ is linear in $t$. In particular every potential
in the class is stable with constant $B_0$.

Each such potential is finite and locally bounded away from the origin
and diverges to $+\infty$ there. Its tail is integrable, since
$\varphi=-S$ outside the unit ball. If one coefficient is positive
and its range is greater than 4, the potential is strictly negative
uniformly on $1\le |x|\le2$: the corresponding rescaled kernel has
a positive minimum there, while the core vanishes. Thus a potential
with this additional simple condition has all the admissibility
properties in the introduction. We next establish its thermodynamics,
uniformly for the larger class \eqref{eq:physical-class}.

\subsection{The canonical limit for every density}

Fix for the moment a value of $t$ satisfying \eqref{eq:t-choice} and
a potential $\varphi$ in \eqref{eq:physical-class}. Write
$Z_{L,n}(\beta)=Z_{\Lambda_L,n}^\varphi(\beta)$, including
$Z_{L,0}=1$, and set
\[
 a_L(\beta,n)=-\frac{\log Z_{L,n}(\beta)}{tL^3}.
\]
We use $0\log0=0$. The following proposition includes convergence
along every sequence of particle numbers with a prescribed limiting
density; that form will be needed in the long-range comparison.

\begin{proposition}[Canonical and grand-canonical limits]
\label{prop:thermodynamics}
For every $\beta>0$ and $d\ge0$ there is a finite limit
\begin{equation}
 H_\varphi(\beta,d)
 =-\lim_{\substack{L\to\infty\\n/L^3\to d}}
       \frac{\log Z_{L,n}(\beta)}{tL^3},
 \qquad H_\varphi(\beta,0)=0.
 \label{eq:canonical-limit}
\end{equation}
Extend $H_\varphi(\beta,d)$ by $+\infty$ for $d<0$. It is closed
and convex in $d$, and finite and concave in $\beta$ for every
$d\ge0$. Moreover,
\begin{equation}
 H_\varphi(\beta,d)
 \ge \frac{\beta\kappa}{t}d^2
       -\frac{\beta B_0}{t}d
       +\frac{d\log d-d}{t}
 \qquad(d\ge0).
 \label{eq:coercivity}
\end{equation}
Let $\Omega=(0,\infty)\times\R$, with coordinates $(\beta,h)$,
and define
\[
 \Xi_L^\varphi(\beta,h)
   =\sum_{n\ge0}e^{thn}Z_{L,n}(\beta).
\]
Then the finite convex pressure is
\begin{equation}
 Q_\varphi(\beta,h)
 =\lim_{L\to\infty}\frac{\log\Xi_L^\varphi(\beta,h)}{tL^3}
 =\sup_{d\ge0}\{hd-H_\varphi(\beta,d)\}.
 \label{eq:pressure-duality}
\end{equation}
The convergence is uniform on compact subsets of $\Omega$, and
\begin{equation}
 H_\varphi(\beta,d)
 =\sup_{h\in\R}\{hd-Q_\varphi(\beta,h)\}
 \qquad(d\in\R).
 \label{eq:partial-conjugacy}
\end{equation}
The bounds above use the same constants for the entire class
\eqref{eq:physical-class} with the chosen $t$.
\end{proposition}

\begin{proof}
We give the details of the canonical limit first. Two finite-volume
bounds control the density adjustments in its construction.

\emph{A trial bound on bounded density intervals.}
Fix $D_0<\infty$ and choose a grid spacing $q\in(0,1/2)$ with
$q^{-3}>D_0+1$. For large $L$, at least $D_0L^3$ grid centers lie
in $\Lambda_L$ at distance $q/8$ from its boundary. For
$n\le D_0L^3$, put one particle in each of $n$ distinct balls of
radius $q/8$ around such centers. Every pair has distance at least
$3q/4$, and each particle has at most $Cq^{-3}$ neighbors at distance
less than one. The core is bounded on $[3q/4,1]$, so the positive
energy of these configurations is at most $E_q n$ for a finite
constant $E_q$. We need not count the negative energy. The $n!$
assignments of the labeled particles to the balls have disjoint
integration regions and cancel the canonical factorial. If $v_q>0$
is the volume of a ball, this gives
\[
 Z_{L,n}(\beta)\ge v_q^n e^{-\beta E_q n}.
\]
Consequently
\begin{equation}
 a_L(\beta,n)\le C_{D_0,\beta}
 \qquad(0\le n\le D_0L^3, L\text{ sufficiently large}).
 \label{eq:trial-bound}
\end{equation}
The constant is uniform over the physical class, and can be chosen
uniformly for $\beta$ in a compact subset of $(0,\infty)$.

\emph{A lower bound.}
There are at most $(L/b+2)^3$ cropped cells in $\Lambda_L$.
Cauchy--Schwarz in \eqref{eq:superstable} and the inequality
$n!\ge(n/e)^n$ therefore give, with $d_L=n/L^3$,
\begin{equation}
 a_L(\beta,n)
 \ge \frac{\beta\kappa}{t(1+2b/L)^3}d_L^2
       -\frac{\beta B_0}{t}d_L
       +\frac{d_L\log d_L-d_L}{t}.
 \label{eq:finite-coercivity}
\end{equation}
Dropping the positive quadratic term also gives the same inequality
with only its last two terms. In particular, $Z_{L,n}$ is positive
and finite, and $a_L$ is bounded above and below when the density
stays in a fixed bounded interval and $L$ is large.

\emph{Convergence at a positive density.}
Fix $\beta>0$ and $d>0$. The increasing infima
\[
 I_j=\inf\left\{a_L(\beta,n):
       L\ge j,\ \left|n/L^3-d\right|\le1/j\right\}
\]
converge to a finite number $I$. Select $\ell_j\ge j$ and $n_j$
in this set such that $a_{\ell_j}(\beta,n_j)\le I_j+1/j$.
Then $n_j/\ell_j^3\to d$ and
$a_{\ell_j}(\beta,n_j)\to I$. Every sequence with
$L\to\infty$ and $n/L^3\to d$ has liminf at least $I_j$ for
each $j$, and hence at least $I$.

To prove the reverse inequality, tile part of a large cube by
\[
 k=\left\lfloor\frac{L}{\ell_j+2}\right\rfloor^3
\]
cubes of side $\ell_j$, leaving gaps of width two. Every interaction
between distinct cubes is nonpositive. For any prescribed counts
$m_1,\ldots,m_k$, restricting the integral to these cubes gives
\begin{equation}
 Z_{L,\sum_i m_i}(\beta)
 \ge\prod_{i=1}^k Z_{\ell_j,m_i}(\beta).
 \label{eq:separated-boxes}
\end{equation}
Indeed there are $(\sum_i m_i)!/\prod_i m_i!$ assignments of labels
to the cubes, and these give exactly the factorials on the right.

Start with $n_j$ particles in every small cube. Define a higher count
$n_{\mathrm{hi}}=\lfloor(4d+4)\ell_j^3\rfloor$ and the diluted
densities
\[
 d_b=\frac{n_j}{(\ell_j+2)^3},\qquad
 d_{\mathrm{hi}}=\frac{n_{\mathrm{hi}}}{(\ell_j+2)^3}.
\]
For large $j$, $n_j>0$, $n_{\mathrm{hi}}>n_j$, and
$d_b\to d$, $d_{\mathrm{hi}}\to4d+4$. If a target count $n$
exceeds $kn_j$, change some cubes to $n_{\mathrm{hi}}$; if it is
smaller, empty some cubes. One additional intermediate count realizes
$n$ exactly. For fixed sufficiently large $j$, this is possible for
all large $L$ with $n/L^3\to d$. The limiting upper fraction of
changed cubes is at most
\[
 c_j=\max\left\{
 \frac{(d-d_b)_+}{d_{\mathrm{hi}}-d_b},
 \frac{(d_b-d)_+}{d_b}\right\},
 \qquad c_j\longrightarrow0.
\]
All new counts are bounded by $(4d+4)\ell_j^3$, so their values of
$a_{\ell_j}$ have a common upper bound by \eqref{eq:trial-bound}.
The baseline values are bounded in absolute value. Applying
\eqref{eq:separated-boxes}, taking the large-volume limit with $j$
fixed, and then letting $j\to\infty$ therefore gives
\[
 \limsup_{\substack{L\to\infty\\n/L^3\to d}}a_L(\beta,n)
 \le\lim_{j\to\infty}
 \left(\frac{\ell_j}{\ell_j+2}\right)^3
 \bigl(a_{\ell_j}(\beta,n_j)+C c_j\bigr)=I.
\]
The single intermediate cube contributes zero in the first limit.
This proves the all-sequence convergence for $d>0$.

\emph{Zero density and convexity.}
When $n/L^3\to0$, the linear lower bound derived above has limit
zero. For an upper bound, fix a sufficiently large side $\ell$ and
use separated cubes as in \eqref{eq:separated-boxes}, with empty
baseline. Fill some cubes with $\lfloor\ell^3\rfloor$ particles
and at most one with an intermediate count. Their number is
$o(L^3)$, so \eqref{eq:trial-bound} gives a limiting upper bound
zero. Thus $H_\varphi(\beta,0)=0$.

For convexity, take $0\le d_1<d_2$ and a density strictly between
them. For a large fixed $\ell$, use counts
$\lfloor d_1\ell^3\rfloor$ and $\lfloor d_2\ell^3\rfloor$
in the separated cubes. The desired density lies strictly between
their diluted densities for all sufficiently large $\ell$.
Choose their proportions to attain the target count, with at most
one intermediate cube. Inequality~\eqref{eq:separated-boxes},
followed first by $L\to\infty$ and then by $\ell\to\infty$,
gives the convexity inequality for $H_\varphi(\beta,\cdot)$.
The endpoint cases of that inequality are immediate.

The function is finite and convex on $[0,\infty)$, hence continuous
on $(0,\infty)$. The linear lower bound gives lower
semicontinuity at zero; since its value there is zero, extension by
$+\infty$ to negative densities is a closed convex function.
For each fixed density, H\"older's inequality makes
$-\log Z_{L,n}(\beta)$ concave in $\beta$. Passage to the limit
gives finite concavity of $H_\varphi(\cdot,d)$ on $(0,\infty)$.
Finally \eqref{eq:finite-coercivity} gives
\eqref{eq:coercivity}.

We have established the canonical limit with enough control over
particle number to pass to the grand-canonical ensemble. In
particular, the density variable in the next maximization will not
require a separate approximation of the canonical free energy.

\emph{Pressure and conjugacy.}
On a compact subset of $\Omega$, \eqref{eq:finite-coercivity}
implies that, for sufficiently large $D$ and all large $L$,
\[
 e^{thn}Z_{L,n}(\beta)\le e^{-c n^2/L^3}
 \qquad(n>DL^3)
\]
with $c>0$ independent of the potential in the class. The sum of
these terms is exponentially small in $L^3$, whereas the term
$n=0$ is one. Among the remaining at most $\lfloor DL^3\rfloor+1$
terms, choose one of maximal size. Every subsequence has a further
subsequence on which its density converges to some $d\in[0,D]$.
The all-sequence canonical limit identifies the limiting exponent
as $hd-H_\varphi(\beta,d)$. Conversely, the term
$n=\lfloor dL^3\rfloor$ supplies this exponent as a lower bound
for every $d\ge0$. This proves \eqref{eq:pressure-duality}.

The finite-volume logarithms are finite and jointly convex in
$(\beta,h)$, by H\"older's inequality applied to their defining
integrals and sums. The limit is finite by the density bound, and
it is nonnegative because the empty configuration is allowed.
Lemma~\ref{lem:elementary-convex-convergence}, the finite-dimensional
convex convergence principle of \cite[Corollary~7.18]{RockafellarWets1998},
upgrades pointwise convergence
to uniform convergence on compact subsets of the open convex set
$\Omega$.
Finally, the closed convexity in density proved above permits the
one-dimensional Fenchel biconjugacy identity, proved in
Lemma~\ref{lem:elementary-conjugacy} and treated in
\cite[Theorem~11.1]{RockafellarWets1998}, yielding
\eqref{eq:partial-conjugacy}.
\end{proof}

\subsection{The quadratic density transform}

We will compare a physical pressure with the pressure obtained by
adding an attraction whose range tends to infinity. The appropriate
variational expression is a quadratic modification in the density
variable. For any physical pressure $Q=Q_\varphi$, write
\[
 H_Q(\beta,d)=\sup_{h\in\R}\{hd-Q(\beta,h)\}.
\]
By \eqref{eq:partial-conjugacy}, this is precisely the canonical
function $H_\varphi$, including its value $+\infty$ for $d<0$.
For $0\le\sigma\le A+1$, define
\begin{equation}
 (T_\sigma Q)(\beta,h)
 =\sup_{d\ge0}
   \left\{hd-H_Q(\beta,d)+\frac{\beta\sigma}{2}d^2\right\}.
 \label{eq:transform}
\end{equation}
This definition does not assert that $T_\sigma Q$ is itself the
pressure of a finite-range pair potential. The next section proves
the approximation by physical pressures that we shall use.

\begin{proposition}[Bounds for the density transform]
\label{prop:transform}
Let $t$ satisfy \eqref{eq:t-choice}. For every physical pressure
$Q_\varphi$ in \eqref{eq:physical-class} and every
$\sigma\in[0,A+1]$, the function $T_\sigma Q_\varphi$ is finite
and convex on $\Omega$, nonnegative, and nondecreasing in $h$.
The supremum in \eqref{eq:transform} is attained. For every compact
$W\subset\Omega$, there is $D_W<\infty$, depending only on
$W,t,b,w,A$, such that every maximizing density belongs to
$[0,D_W]$, uniformly over all potentials in the class and all these
values of $\sigma$. The transformed pressures are also uniformly
bounded on $W$.
\end{proposition}

\begin{proof}
The value at $d=0$ is zero. By \eqref{eq:coercivity}, the objective
in \eqref{eq:transform} is bounded above by
\begin{equation}
 hd-\beta\left(\frac\kappa t-\frac\sigma2\right)d^2
       +\frac{\beta B_0}{t}d
       -\frac{d\log d-d}{t}.
 \label{eq:transform-coercivity}
\end{equation}
Condition~\eqref{eq:t-choice} gives a uniform positive lower bound
on $\kappa/t-\sigma/2$. On a compact set of $(\beta,h)$ with
$\beta$ bounded away from zero, the right side is therefore
negative for every sufficiently large $d$, uniformly in the
potential and in $\sigma$. This proves the common density bound
and the finite upper bound for the supremum. The objective is upper
semicontinuous in $d$, since $H_Q(\beta,\cdot)$ is lower
semicontinuous. It attains its maximum on the resulting compact
interval.

For each fixed $d\ge0$, the function in braces in
\eqref{eq:transform} is convex in $(\beta,h)$: the canonical
function is concave in $\beta$, and the other terms are affine.
Their supremum is convex. Its monotonicity in $h$ follows from
$d\ge0$, and its nonnegativity from the zero-density choice.
\end{proof}

The normalizations are now fixed. In particular, for every positive
density the free energy in the statement of the problem is
\begin{equation}
 f_\varphi(\beta,d)=\frac{t}{\beta}H_\varphi(\beta,d).
 \label{eq:free-energy-normalization}
\end{equation}
For a real-valued function $F$ on a compact set $K$, write
$\|F\|_K=\sup_{x\in K}|F(x)|$.
All limits and finiteness statements in this section hold for every
$\beta>0$, including for infinite sums in
\eqref{eq:physical-class}. The remaining task is to choose the
attractive ranges and coefficients so that one of these canonical
functions has two different temperature slopes at one fixed
positive density.
